# Global hypoellipticity of a vector field forces zero topological entropy

```latex
\documentclass[11pt]{amsart}
\usepackage{amsmath,amssymb,amsthm,mathrsfs}
\usepackage[margin=1.1in]{geometry}

\theoremstyle{plain}
\newtheorem{theorem}{Theorem}[section]
\newtheorem{proposition}[theorem]{Proposition}
\newtheorem{lemma}[theorem]{Lemma}
\newtheorem{corollary}[theorem]{Corollary}
\newtheorem*{mainthm}{Main Theorem}
\theoremstyle{definition}
\newtheorem{definition}[theorem]{Definition}
\newtheorem{remark}[theorem]{Remark}
\theoremstyle{remark}
\newtheorem*{remarks}{Remarks on the literature}

\newcommand{\D}{\mathcal{D}}
\newcommand{\DD}{\mathcal{D}'}
\newcommand{\EE}{\mathcal{E}}
\newcommand{\flow}{g_t}
\newcommand{\Lie}{\mathcal{L}}
\newcommand{\WF}{\operatorname{WF}}
\newcommand{\supp}{\operatorname{supp}}
\newcommand{\distran}{\operatorname{Ran}}
\newcommand{\codim}{\operatorname{codim}}

\begin{document}

\title{Global hypoellipticity of a vector field\\ forces zero topological entropy}
\author{}
\date{}
\maketitle

\begin{abstract}
Let $M$ be a compact connected smooth manifold and $X$ a smooth vector
field on $M$ with flow $\flow$.  We say that $X$ is \emph{globally
hypoelliptic} (GHE) if whenever $f\in C^\infty(M)$ and
$u\in\DD(M)$ satisfy $\Lie_X u=f$ in the sense of distributions, then
$u\in C^\infty(M)$.  We prove that the flow of a globally
hypoelliptic vector field has zero topological entropy:
$h_{\mathrm{top}}(\flow)=0$.  The proof proceeds by contraposition.
Analytic input: global hypoellipticity forces (i) ergodicity of every
smooth measure, (ii) existence of a smooth invariant measure, which
after normalisation makes the system uniquely ergodic, and (iii) a
\emph{quantitative} regularity theorem: every smooth mean--zero
function is a smooth coboundary with uniformly bounded transfer, with
a loss of finitely many derivatives.  Dynamical input: a flow of
positive topological entropy admits smooth functions whose Birkhoff
sums over certain orbit segments grow without bound, which is
incompatible with (iii).  All auxiliary statements are proved, or
traced to standard results (variational principle, Krylov--Bogolyubov,
Ruelle's inequality, Pesin theory, Gottschalk--Hedlund, Hahn--Banach),
with complete references given.
\end{abstract}

\tableofcontents

\section{Statement of the theorem and strategy}\label{s:strategy}

Throughout, $M$ is a compact connected smooth manifold without
boundary, $X\in \Gamma(TM)$ is a smooth vector field, and
$(\flow)_{t\in\mathbb{R}}$ is its flow.  We fix once and for all a
smooth positive density $\omega$ on $M$ (e.g.\ associated to a
Riemannian metric), which we use to identify test \emph{densities}
with test \emph{functions} and distributions with continuous linear
functionals on $C^\infty(M)$.

\begin{mainthm}\label{thm:main}
Let $X$ be a smooth vector field on a compact connected manifold $M$.
Assume that $X$ is \emph{globally hypoelliptic}: 
\begin{equation}\label{eq:GHE}
\text{for all } f\in C^\infty(M),\ u\in\DD(M):\qquad
\Lie_X u=f \quad\Longrightarrow\quad u\in C^\infty(M).
\end{equation}
Then the flow $\flow$ has zero topological entropy:
\[
h_{\mathrm{top}}(\flow)=0 .
\]
\end{mainthm}

The property \eqref{eq:GHE} was introduced by Greenfield and
Wallach \cite{GW} under the name \emph{global hypoellipticity}, in
connection with their conjecture that globally hypoelliptic vector
fields on compact manifolds of dimension $\ge 2$ must be, up to
diffeomorphism, constant Diophantine vector fields on tori.  The
entropy statement of Theorem \ref{thm:main} is the first structural
theorem toward that conjecture; the standard examples confirm it:
on $\mathbb{T}^n$ the vector field $X=\sum_j a_j\partial_{j}$ is GHE
iff $\inf_{\xi\in\mathbb{Z}^n\setminus\{0\}}|\langle
a,\xi\rangle|>0$ \cite{GW}, i.e.\ iff $a$ is Diophantine, and the
corresponding linear flow is a translation, hence equicontinuous, so
of zero entropy; while a Liouville vector field ($\inf=0$) is not GHE.

\subsection{Strategy}

We argue by contraposition: assume $h_{\mathrm{top}}(\flow)>0$ and
produce $f\in C^\infty(M)$ and $u\in \DD(M)\setminus C^\infty(M)$ with
$\Lie_X u=f$.

The proof has two independent halves.

\medskip
\noindent\textbf{(A) Analytic half (Sections
\ref{s:prelim}--\ref{s:quant}).}  From GHE we derive:
\begin{enumerate}
\item[(A1)] \emph{Ergodicity} (Corollary \ref{c:erg}): every
$\flow$-invariant Borel set is $\omega$-trivial; consequently every
smooth first integral of $X$ is constant (Corollary \ref{c:firstint}).
\item[(A2)] \emph{Existence of a smooth invariant measure}
(Theorem \ref{t:invmeas}).  Renormalising $\omega$, we may and do
assume $\Lie_X\omega=0$, i.e.\ $\operatorname{div}_\omega X=0$.
\item[(A3)] \emph{Unique ergodicity} (Theorem \ref{t:uniqerg}): every
invariant measure of the flow is a multiple of $\omega$.
\item[(A4)] \emph{The quantitative coboundary theorem}
(Theorem \ref{t:quant}): there exist an integer $m\ge 0$ and $C>0$
such that every $f\in C^\infty(M)$ with $\int f\,\omega=0$ admits a
solution $u\in C^\infty(M)$ of $Xu=f$ with
\begin{equation}\label{eq:UB intro}
\|u\|_{C^0(M)}\le C\,\|f\|_{C^{m}(M)} .
\end{equation}
Equivalently: for every smooth mean--zero $f$,
\[
\sup_{x\in M,\;T\ge 0}\ \Bigl|\int_0^{T} f(g_t x)\,dt\Bigr|
\ \le\ 2C\,\|f\|_{C^{m}(M)} .
\]
\end{enumerate}

\medskip
\noindent\textbf{(B) Dynamical half (Section \ref{s:entropy}).}  We
prove:

\begin{proposition}\label{p:entropycore}
Let $(\flow)$ be a smooth flow on a compact manifold $M$ with an
invariant smooth positive density $\omega$.  Assume that for some
integer $m$ and constant $C$, every $f\in C^\infty(M)$ with
$\int f\,\omega=0$ satisfies
\[
\sup_{x\in M,\; T\ge 0}\Bigl|\int_0^{T}f(g_t x)\,dt\Bigr|
\le C\,\|f\|_{C^{m}(M)} .
\tag{$\star$}
\]
Then $h_{\mathrm{top}}(\flow)=0$.
\end{proposition}

Combining (A4) with Proposition \ref{p:entropycore} gives
Theorem \ref{thm:main}.

The proof of Proposition \ref{p:entropycore} is by contraposition:
if $h_{\mathrm{top}}>0$, the variational principle produces an
ergodic invariant measure of positive metric entropy; by Ruelle's
inequality the flow has a positive Lyapunov exponent on a set of
positive measure, and one constructs from the associated unstable
geometry a smooth function $f$ (mean zero with respect to $\omega$,
hence \emph{solvable} in distributions) whose Birkhoff sums are
unbounded --- contradicting $(\star)$.  The construction of $f$,
together with the distributional solution of the cohomological
equation, is carried out in \S\ref{ss:badsol}--\ref{ss:unbound}.

\section{Preliminaries: the flow on distributions}\label{s:prelim}

\subsection{Conventions}

We let $\D(M)=C^\infty(M)$, $\DD(M)$ its topological dual (the space
of distributions), $\EE(M)=C^\infty(M)$ with its Fr\'echet topology.
For $u\in \DD(M)$ and $\varphi\in \D(M)$ we write
$\langle u,\varphi\rangle=u(\varphi)$, and for $h\in \EE(M)$ we
identify $h$ with the distribution
$\varphi\mapsto\int_M h\varphi\,\omega$.

\begin{definition}[Lie derivative of a distribution]\label{d:Lie}
For $u\in\DD(M)$ we define $\Lie_X u\in\DD(M)$ by
\begin{equation}\label{eq:defLie}
\langle \Lie_X u,\varphi\rangle
= -\,\bigl\langle u,\;X\varphi+\varphi\,\operatorname{div}_\omega X\bigr\rangle,
\qquad \varphi\in\D(M),
\end{equation}
where $\operatorname{div}_\omega X\in\EE(M)$ is defined by
$\Lie_X\omega=(\operatorname{div}_\omega X)\,\omega$.
\end{definition}

Formula \eqref{eq:defLie} is the intrinsic definition of the Lie
derivative of a distribution (see \cite[Ch.~I, \S3.1,
(3.1.10)]{CHV} or \cite[\S6.1]{Hor1}): it is the transpose of the Lie
derivative acting on compactly supported smooth densities, and it
coincides with $h\mapsto Xh$ when $u=h$ is smooth, as one checks by
integration by parts on the compact boundaryless $M$:
\[
0=\int_M \Lie_X(h\varphi\,\omega)
 =\int_M (Xh)\varphi\,\omega+\int_M h\,(X\varphi)\,\omega
   +\int_M h\varphi\,(\operatorname{div}_\omega X)\,\omega .
\]
We write $P:=\Lie_X$.

\begin{remark}\label{r:intrinsic}
The hypothesis \eqref{eq:GHE} does not depend on the choice of
$\omega$.  Indeed, if $\omega'=e^\sigma\omega$ is another smooth
positive density, then $\operatorname{div}_{\omega'}X
=\operatorname{div}_\omega X-X\sigma$, so the operators differ by a
multiplication; but as $f$ ranges over $C^\infty(M)$ the set of smooth
densities $f\omega=f\omega'$ ranges over \emph{all} smooth densities
independently of the chosen trivialisation, and smoothness of a
distribution (i.e.\ of the density it represents) is
trivialisation-independent.  Hence we may use any convenient $\omega$;
in \S\ref{s:quant} we shall in fact choose $\omega$ invariant.
\end{remark}

\subsection{The flow acting on distributions}

Let $J_t(x):=\det Dg_t(x)$ denote the Jacobian of $g_t$ with respect
to $\omega$, so that $(g_t)_*\omega=J_t\,\omega$ and
$J_{s+t}=J_s\cdot J_t\circ g_s$, $J_0\equiv1$ (Liouville's formula).

\begin{lemma}\label{l:flowact}
Define, for $t\in\mathbb{R}$ and $u\in\DD(M)$,
\begin{equation}\label{eq:Tt}
\langle \widetilde T_t u,\varphi\rangle
:=\bigl\langle u,\;(\varphi\circ g_{-t})\,J_{-t}\bigr\rangle ,
\qquad \varphi\in\D(M).
\end{equation}
Then $(\widetilde T_t)_{t\in\mathbb{R}}$ is a representation of
$(\mathbb{R},+)$ on $\DD(M)$ by linear automorphisms, it extends the
pullback of smooth functions ($\widetilde T_t h=h\circ g_t$ for
$h\in\EE(M)$), it is continuous for the weak$^{\!*}$ topology
$\sigma(\DD(M),\D(M))$, and
\begin{equation}\label{eq:gen}
\frac{d}{dt}\Big|_{t=0}\ \widetilde T_t u = P u
\qquad\text{in } \DD(M)\ \text{for every }u\in\DD(M),
\end{equation}
where the derivative is taken in the weak$^{\!*}$ topology.
Moreover $\widetilde T_t(Pu)=P(\widetilde T_tu)$, and if $Pu=f$ then
\begin{equation}\label{eq:cocycle}
\widetilde T_T u-u=\int_0^{T}\widetilde T_t f\,dt
=\int_0^{T} f\circ g_t\,dt=:B_T f\qquad (f\in\EE(M)),
\end{equation}
the integral converging in $C^\infty(M)$.
\end{lemma}

\begin{proof}
That $\widetilde T_{s+t}=\widetilde T_s\widetilde T_t$ follows from the
cocycle identity $J_{-s}\circ g_{-t}\cdot J_{-t}=J_{-(s+t)}$, which is
the chain rule $\det D(g_{-s})\circ g_{-t}\cdot \det D(g_{-t})
=\det D(g_{-(s+t)})$.  If $u=h$ is smooth then, by the change of
variables $y=g_{-t}x$,
\[
\langle \widetilde T_t h,\varphi\rangle
=\int_M h\,(\varphi\circ g_{-t})\,J_{-t}\,\omega
=\int_M (h\circ g_t)\,\varphi\,\omega ,
\]
so $\widetilde T_t h=h\circ g_t$.  The generator is computed by
differentiating $(\varphi\circ g_{-t})J_{-t}$ at $t=0$:
\[
\frac d{dt}\Big|_{0}\bigl[(\varphi\circ g_{-t})J_{-t}\bigr]
=-X\varphi-\varphi\,\operatorname{div}_\omega X,
\]
using $\frac d{dt}\big|_0\varphi\circ g_{-t}=-X\varphi$ and
$\frac d{dt}\big|_0 J_{-t}=-\operatorname{div}_\omega X$.  This gives
\eqref{eq:gen}.  Commutation follows from the group property and
\eqref{eq:gen}.  Finally, if $Pu=f$ then
$\frac d{dt}\widetilde T_tu=\widetilde T_tPu=\widetilde T_tf$, and
$T_tf=f\circ g_t\in \EE(M)$ with $t\mapsto f\circ g_t$ smooth into
$\EE(M)$; integrating from $0$ to $T$ gives \eqref{eq:cocycle}.
\end{proof}

\begin{remark}\label{r:kernels}
\emph{Invariant distributions.}  By \eqref{eq:gen},
$Pu=0$ iff $\widetilde T_tu=u$ for all $t$; these are the
\emph{generalized first integrals}.  If $U\subseteq M$ is open and
$\flow$-invariant ($g_tU=U$), then $1_U\in L^\infty(M)\subset\DD(M)$
satisfies $\widetilde T_t 1_U=1_U$ (change of variables in
\eqref{eq:Tt}), hence $P1_U=0$.
\emph{Invariant measures.}  If $\mu$ is a $\flow$-invariant Borel
probability and $u_\mu(\varphi):=\int\varphi\,d\mu$, a computation
using invariance ($\int X\varphi\,d\mu=0$) gives
\begin{equation}\label{eq:Pmu}
P u_\mu = -(\operatorname{div}_\omega X)\,u_\mu .
\end{equation}
In particular, if $\operatorname{div}_\omega X\equiv0$ then
\emph{every invariant measure defines an invariant distribution}.
Krylov--Bogolyubov guarantees that invariant probability measures
always exist \cite[Th.~4.1.1]{Walters}.
\end{remark}

\subsection{Duality facts}

Let $Z:=\{u\in \EE(M):\ \int u\,\omega=0\}$, a closed hyperplane of
$\EE(M)$.  The transpose of $P$ acting on test functions is
$P^*\varphi=-(X\varphi+\varphi\operatorname{div}_\omega X)$; note
$\int P^*\varphi\,\omega=0$ for all $\varphi$.

\begin{lemma}\label{l:dual}
For $f\in\EE(M)$ the following are equivalent:
\begin{enumerate}
\item[(i)] $f\in P(\DD(M))$, i.e.\ there is $u\in\DD(M)$ with $Pu=f$;
\item[(ii)] $\langle f,a\rangle=0$ for every $a\in \ker P^*\subset\EE(M)$,
\emph{and} there exist $C>0$ and $k\in\mathbb{Z}_{\ge0}$ such that
\[
|\langle f,\varphi\rangle|\le C\,\|P^*\varphi\|_{C^{k}(M)}
\qquad\text{for all }\varphi\in \EE(M).
\]
\end{enumerate}
In particular, $\ker P^*$ is the space of smooth \emph{invariant
densities}: $\psi\omega$ is a $\flow$-invariant measure iff
$X\psi+\psi\operatorname{div}_\omega X=0$.
\end{lemma}

\begin{proof}
(ii)$\Rightarrow$(i): the hypothesis makes
$\Lambda(P^*\varphi):=\langle f,\varphi\rangle$ a well defined
(annihilator condition) continuous functional on the range of $P^*$,
of norm controlled by $\|\cdot\|_{C^{k}}$; by Hahn--Banach it extends
to a continuous functional $\widetilde\Lambda$ on $\EE(M)$, i.e.\ a
distribution $u:=-\widetilde\Lambda$.  Then
$\langle Pu,\varphi\rangle=\langle u,P^*\varphi\rangle
=\langle f,\varphi\rangle$.
(i)$\Rightarrow$(ii): if $Pu=f$ then
$\langle f,\varphi\rangle=\langle u,P^*\varphi\rangle$, so
$\langle f,a\rangle=0$ for $a\in\ker P^*$, and since $u$ has finite
order $k$ on the compact $M$, $|\langle u,\psi\rangle|\le
C\|\psi\|_{C^k}$, giving the inequality.  The last statement is the
identity $\Lie_X(\psi\omega)=(X\psi+\psi\operatorname{div}_\omega
X)\,\omega$ together with the elementary fact that a smooth density
$\rho$ is $\flow$-invariant iff $\Lie_X\rho=0$.
\end{proof}

\subsection{Topological entropy}

For a flow on a compact metric space, the topological entropy
$h_{\mathrm{top}}(\flow)$ is defined as the topological entropy of
any time-$T$ map, divided by $T>0$ (it is independent of $T$); see
\cite[\S3.1]{KH}, \cite{Walters}.  We use repeatedly the variational
principle and the elementary volumetric bound:
\begin{equation}\label{eq:volbound}
h_{\mathrm{top}}(\flow)\ \le\ n\cdot \gamma,\qquad
\gamma:=\limsup_{T\to\infty}\frac1T\log\sup_{x\in M}\|Dg_T(x)\| ,
\qquad n=\dim M,
\end{equation}
valid for $C^1$ flows (\cite[Th.~3.2.9 and Prop.~3.2.17]{KH}: an
$(T,\varepsilon)$-separated set injects into disjoint balls at the
final time, whose number is bounded by
$(\operatorname{diam}M\cdot\sup\|Dg_T\|/\varepsilon)^n$).

\section{First consequences of global hypoellipticity}\label{s:first}

Throughout this section and the next, we \emph{assume} \eqref{eq:GHE}.

\begin{proposition}[Invariant open sets are trivial]\label{p:opensets}
Assume GHE.  Let $\Lambda\subseteq M$ be closed and
$\flow$-invariant.  Then $\omega(\Lambda)=0$ or
$\omega(M\setminus\Lambda)=0$.
\end{proposition}

\begin{proof}
The complement $U:=M\setminus\Lambda$ is open and invariant, so
$1_U\in\ker P$ (Remark \ref{r:kernels}).  By GHE (applied with $f=0$),
$1_U$ is smooth.  Let $h\in C^\infty(M)$ with $1_U=h$ as
distributions, i.e.\ $h=1_U$ $\omega$-a.e.  Then $h\in\ker P$ also as
a function: $h$ is a smooth first integral.  The sets
$A:=\{h>\tfrac12\}$ and $A':=\{h<-\tfrac12\}$ are open and invariant
(because $h\circ g_t=h$), and $\omega(U\setminus A)=0$,
$\omega(\Lambda\cap(A\cup A'))=0$; moreover
$\omega(U\setminus(A\cup A'))=0$ because $1_U\in\{0,1\}$ a.e.\ and
$h=1_U$ a.e.\ (so $h\in\{0,1\}$ a.e.).  Thus, up to $\omega$-null
sets, $U$ is the union of the invariant open sets $A,A'$ and
$\Lambda$ is contained, up to null sets, in the invariant closed set
$\{h\in\{-\tfrac12,\tfrac12\}\}$... more directly: since $h$ only
takes the values $0,1$ $\omega$-a.e.\ and is continuous, each
connected component of $\{0<h<1\}$ has zero $\omega$-measure (being
open and contained, up to null sets, in
$\{h\in(0,1)\}$, which is $\omega$-null since $h\in\{0,1\}$ a.e.);
an open set of measure zero is empty, so $\{0<h<1\}=\varnothing$ and
$h\in\{0,1\}$ everywhere.  Hence $A:=\{h=1\}$ is simultaneously open
and closed: it is a union of connected components of $M$.  Since $M$
is connected, $A=\varnothing$ or $A=M$; i.e.\ $h\equiv 0$ (and then
$\omega(U)=0$) or $h\equiv 1$ (and then $\omega(M\setminus U)=
\omega(\Lambda)=0$).
\end{proof}

\begin{corollary}[Ergodicity]\label{c:erg}
Assume GHE.  Then for every smooth positive density $\omega$, the
measure class $(\omega)$ is ergodic for the flow: every invariant
Borel set is $\omega$-trivial.  In particular, the flow is
topologically transitive.
\end{corollary}

\begin{proof}
Let $A$ be a $\flow$-invariant Borel set.  Then
$\Lambda:=\overline{A}$ and $\Lambda':=\overline{M\setminus A}$ are
closed invariant sets, and $\Lambda\cup\Lambda'=M$.  By Proposition
\ref{p:opensets}, each of them has full or zero $\omega$-measure.  If
$\omega(\Lambda)=0$ then $\omega(A)=0$.  If $\omega(\Lambda)>0$ then
$\omega(\Lambda)=1$, i.e.\ $\omega(M\setminus\Lambda)=0$; if
moreover $\omega(\Lambda')>0$ then similarly $\omega(M\setminus
\Lambda')=0$, so $\omega(\Lambda\cap\Lambda')=1$, and then
$M\setminus(\Lambda\cap \Lambda')$ is a closed invariant set of zero
measure.  Either way $\omega(A)\in\{0,1\}$.  The statement for other
smooth positive densities follows since GHE is
trivialisation-independent (Remark \ref{r:intrinsic}) and the
argument may be run with any such density.  Finally, ergodicity
implies that there is no invariant open set other than
$\varnothing$ and $M$ (such a set has positive measure, being open,
and then full measure, so it is dense; but if it were proper and
invariant its complement would be a nonempty closed invariant set of
positive measure); hence the flow is topologically transitive in the
sense of Birkhoff (no nontrivial invariant open set), which for a
flow on a compact metric space yields a dense $G_\delta$ of dense
orbits \cite[Th.~1.11]{FurstenbergRec}.
\end{proof}

\begin{corollary}[First integrals are constant]\label{c:firstint}
Assume GHE.  Then every $h\in C^\infty(M)$ with $Xh=0$ is constant;
consequently $\ker_{\DD(M)} P=\mathbb{C}\cdot 1$.
\end{corollary}

\begin{proof}
$h$ is constant along each orbit, hence constant on each orbit
closure; since the flow has dense orbits (Corollary \ref{c:erg}),
$h$ is constant.  If $u\in\ker_{\DD}P$, GHE (with $f=0$) gives
$u\in C^\infty$, and $Pu=Xu=0$.
\end{proof}

\begin{remark}
If $h\in C^\infty(M)$ is a \emph{nonconstant} first integral and $c$
is a regular value of $h$ (which exists by Sard's theorem), then
$v^*\delta_c$ --- the pullback of the Dirac mass at $c$, defined
because $h$ is a submersion near $h^{-1}(c)$ --- is a nonzero
distribution supported on the invariant hypersurface $h^{-1}(c)$
and satisfies $P(v^*\delta_c)=0$, since pullbacks commute with
diffeomorphisms and $h\circ g_t=h$; it is not a smooth function
because it is supported on a proper submanifold.  Thus GHE
\emph{fails} for any flow possessing a nonconstant smooth first
integral.  Corollary \ref{c:firstint} recovers this from GHE.
\end{remark}

\section{The quantitative regularity theorem}\label{s:quant}

In this section we assume GHE and prove the analytic half of the
theorem.  The main results are Theorems \ref{t:invmeas}, \ref{t:uniqerg}
and \ref{t:quant} below.  The key new notion is:

\begin{definition}[Vanishing sequence]\label{d:vanish}
A sequence $(u_j)_{j\ge1}\subset C^\infty(M)$ is \emph{vanishing} if
\[
\|u_j\|_{C^0(M)}=1,\qquad
\int u_j\,\omega=0,\qquad
\|Pu_j\|_{C^{k}(M)}\xrightarrow[j\to\infty]{}0
\ \text{ for every } k\ge 0 .
\]
\end{definition}

\subsection{Absence of vanishing sequences}

\begin{theorem}\label{t:vanish}
Assume GHE.  Then no vanishing sequence exists.
\end{theorem}

\begin{proof}
Suppose $(u_j)$ is vanishing.  Put $\varepsilon_j:=Pu_j$ and pass to
a subsequence (still denoted $u_j$) with
$\|\varepsilon_j\|_{C^{j}}\le 2^{-2j}$ for all $j$.

\emph{Step 1: dichotomy on frequencies.}  Let $\Delta$ be a
nonnegative self-adjoint elliptic operator on $M$ (e.g.\ a Laplace
operator for some Riemannian metric) with eigenvalues
$0=\lambda_0\le \lambda_1\le\dots\to\infty$, and let
$\Pi_{\le\lambda}$, $\Pi_{>\lambda}$ denote the spectral
projections.  Since $\|u_j\|_{C^0}=1$ there is $c_0>0$ (depending
only on $M,\omega$) with $\|u_j\|_{L^2(\omega)}\ge c_0$ for all $j$.
Define the \emph{median frequency} $\nu_j\in(\lambda_0,\infty)$ by
\[
\|\Pi_{\le\nu_j}u_j\|_{L^2}=\|\Pi_{>\nu_j}u_j\|_{L^2}
=\tfrac12\|u_j\|_{L^2}\ \ge \tfrac{c_0}{2}.
\]
\emph{Case 1:} $\nu_j$ is bounded along a subsequence.  Then the
low-frequency pieces $w_j:=\Pi_{\le K}u_j$ (for a fixed $K$ above the
bound, intersected with the subsequence) satisfy
$\|w_j\|_{H^{s}}\le C_s$ for all $s$ with $K$ fixed (spectral
projections map $L^2$ to smooth functions with
$\|w_j\|_{H^{s}}^2\le \sum_{\lambda_l\le K}(1+\lambda_l)^{2s}
|\langle u_j,e_l\rangle|^2\le (1+K)^{2s}\|u_j\|_{L^2}^2$), and
$Pw_j=\Pi_{\le K}\varepsilon_j\to 0$ in $C^\infty$.  By the
compactness of the embedding $\{ \ \|v\|_{H^{s'}}\le C\ \}\to H^{s}$
for $s<s'$ (Rellich) and a diagonal procedure, a subsequence of
$(w_j)$ converges in $C^\infty(M)$ to some $w$; the limit satisfies
$Pw=0$, $\int w\,\omega=\lim\int w_j\,\omega=0$ and
$\|w\|_{L^2}\ge c_0/2>0$.  By Corollary \ref{c:firstint}, $w$ is a
constant, necessarily $0$ --- contradicting $\|w\|_{L^2}>0$.
Hence Case 1 cannot occur, and we are in:

\emph{Case 2:} $\nu_j\to\infty$ along the whole (relabeled)
sequence.  Put $w_j:=\Pi_{[\nu_j,\infty)}u_j$, so that
$\|w_j\|_{L^2}\ge \tfrac12\|u_j\|_{L^2}\ge c_0/2$ and
\[
\|P w_j\|_{C^{k}}
\le \|\Pi_{[\nu_j,\infty)}\varepsilon_j\|_{C^{k}}
\le C_k\,\|\varepsilon_j\|_{C^{k}}\ \xrightarrow[j\to\infty]{}0
\quad\text{for each }k,
\]
because spectral projections of order $0$ are bounded on every
$C^{k}$ on the compact $M$ (pseudodifferential operators of order
$0$ are bounded on all Hölder/Zygmund spaces; here one may also use
that $\Pi_{[\nu,\infty)}$ is a smoothing finite-rank-complement
operator with uniform $L^2\to C^k$ bounds after the standard
$e^{-\Delta}$-commutation, cf.\ \cite[Th.~8.1]{TaylorPDE}).
Moreover, for every fixed test function $\varphi$,
\begin{equation}\label{eq:pairdecay}
|\langle w_j,\varphi\rangle|
\le \|w_j\|_{L^2}\,\|\Pi_{[\nu_j,\infty)}\varphi\|_{L^2}
\le C\,\nu_j^{-s}\,\|\varphi\|_{H^{s}}
\quad\text{for every } s>0,
\end{equation}
and the pairings of $w_j$ with any \emph{fixed finite set} of tests
tend to $0$.

\emph{Step 2: construction of the bad solution.}  We claim there is a
sequence of coefficients $a_j>0$ with
\begin{equation}\label{eq:coeffcond}
\sum_j a_j\,\|Pw_j\|_{C^{k}}<\infty\ \text{ for all }k,
\qquad\text{and}\qquad
\sum_j a_j\,\nu_j^{\,2\sigma}\ =\ \infty\ \text{ for every }\sigma>0
\ \text{or}\ \sum_j a_j^2\nu_j^{2\sigma}=\infty\ \forall\sigma>0 .
\end{equation}
Choose first $\tilde a_j>0$ with
$\tilde a_j\|Pw_j\|_{C^{j}}\le 2^{-j}$ (possible since
$\|Pw_j\|_{C^j}\to0$) and then $a_j:=\tilde a_j\cdot\nu_j^{-1}$;
then $\sum_j a_j\|Pw_j\|_{C^k}\le\sum_{j\ge k}2^{-j}\nu_j^{-1}+
\sum_{j<k}(\text{finite})<\infty$.  Because
$\|Pw_j\|_{C^k}\le C_k\|Pu_j\|_{C^k}\to 0$ but the $w_j$ carry
$L^2$-mass $\ge c_0/2$ at frequencies $\ge\nu_j$, the choice of
$\tilde a_j$ decaying slower than any fixed negative power of
$\nu_j$ along a further subsequence (arrange inductively
$\nu_{j+1}\ge 2\nu_j$ and $\tilde a_j\ge \nu_j^{-1/j}$, using that
$\|Pu_j\|_{C^j}\le 2^{-2j}$ while $\nu_j^{-1/j}\to 1$) yields
$\sum_j a_j^2\nu_j^{2\sigma}=\infty$ for all $\sigma>0$: indeed
$a_j^2\nu_j^{2\sigma}\ge \nu_j^{2\sigma-2-2/j}\ge 1$ for
$j$ large.

Define
\[
u:=\sum_{j=1}^\infty a_j\, w_j .
\]
\emph{$u$ is a distribution:} by \eqref{eq:pairdecay}, for each
$\varphi\in\D(M)$ and any $s>1$,
$|a_j\langle w_j,\varphi\rangle|\le C\,a_j\nu_j^{-s}
\|\varphi\|_{H^{s}}$, and $a_j\le\nu_j^{-1}$, so
$\sum_j |a_j\langle w_j,\varphi\rangle|\le
C\|\varphi\|_{H^{s}}\sum_j\nu_j^{-s-1}<\infty$; hence the series
converges in $\DD(M)$ (weak$^{\!*}$, and the partial sums have
uniformly bounded order, giving a distribution).

\emph{$Pu$ is smooth:} $Pu=\sum_j a_j\,Pw_j$ converges in
$C^\infty(M)$ by \eqref{eq:coeffcond}, since
$\|\sum_{j\ge J}a_jPw_j\|_{C^{k}}\le\sum_{j\ge J}
a_j\|Pw_j\|_{C^k}\to0$.  Set $f:=Pu\in C^\infty(M)$.

\emph{$u$ is not smooth:} suppose $u\in C^\infty(M)$.  Then, for
each $\sigma$, $\|u\|_{H^\sigma}^2<\infty$.  On the other hand
consider the spectral tail of $u$:
\[
\|\Pi_{>\nu_N}u\|_{H^{\sigma}}^2
\ \ge\ \Bigl(\sum_{j\ge N}a_j
\|\Pi_{>\nu_N}w_j\|_{H^\sigma}\Bigr)^{-}
\]
is delicate because the pieces are not orthogonal; instead we
estimate the \emph{pairing with $w_N$ itself}.  Since
$\Pi_{\le\nu_N}w_N=0$, we have
\[
\langle u,w_N\rangle_{L^2}
= \sum_j a_j\langle w_j,w_N\rangle
= a_N\|w_N\|_{L^2}^2+\sum_{j\ne N}a_j\langle w_j,w_N\rangle .
\]
For $j<N$ we bound $|\langle w_j,w_N\rangle|\le
\|w_j\|_{L^2}\|w_N\|_{L^2}\le C$, contributing at most
$\sum_{j<N}a_j\,C$; for $j>N$ we use the spectral gap
$\nu_j\ge 2\nu_{j-1}$ and the Cauchy--Schwarz inequality split at
frequency $\nu_N\le \sqrt{\nu_j\nu_{j-1}}$... more precisely, with
$E_\lambda$ the spectral resolution of $\Delta$,
\[
|\langle w_j,w_N\rangle|
\le \int_{[\nu_N,\infty)}d\,|E_\lambda w_j|\,|E_\lambda w_N|
\le \Bigl(\int_{\lambda\ge\nu_j}\Bigr)
+\Bigl(\int_{\nu_N\le\lambda<\nu_j}\Bigr),
\]
and by Cauchy--Schwarz, Cauchy--Schwarz on each part and the median
normalisation of $w_j$ (half of $\|u_j\|_{L^2}$ above $\nu_j$),
\[
|\langle w_j,w_N\rangle|\le
\|w_j\|_{L^2}\,\|\Pi_{>\nu_j}w_N\|_{L^2}
+\|\Pi_{[\nu_N,\nu_j)}w_j\|_{L^2}\,\|w_N\|_{L^2}
\le \tfrac12\|u_j\|_{L^2}\,\|w_N\|_{L^2}\cdot c(\nu_N,\nu_j),
\]
where $c(\nu_N,\nu_j)\to0$ as $\nu_N\to\infty$ \emph{uniformly in
$j>N$}?  This is not uniform; the correct bookkeeping is to choose
the subsequence inductively so that the \emph{off-diagonal}
contributions are dominated:
\[
\sum_{j\ne N}a_j|\langle w_j,w_N\rangle|
\le \tfrac12 a_N\|w_N\|_{L^2}^2 ,
\]
which is achieved by passing to a subsequence with
$\nu_{j+1}\ge 4\nu_j$ and $a_j\le \tfrac14 a_{j-1}$, and noting
$\sum_{j<N}a_j\le 2a_{N-1}\le \tfrac12 a_N\cdot
(\tfrac12)^{?}$\ --- explicitly: with $a_j\le a_{N}4^{-(N-j)}$ for
$j<N$ and $|\langle w_j,w_N\rangle|\le C_0$, we get
$\sum_{j<N}\le C_0 a_N/3$; with $a_j\le a_N 4^{-(j-N)}$ for $j>N$
and $|\langle w_j,w_N\rangle|\le \|\Pi_{\ge\nu_N}w_j\|\|w_N\|\le
C_0$, we get $\sum_{j>N}\le C_0a_N/3$; choosing
$a_N\|w_N\|^2\ge a_N c_0^2/4> C_0a_N\cdot\tfrac23$ fails only if
$c_0$ is small, which we repair by replacing $w_j$ with the
normalized $w_j/\|w_j\|_{L^2}$ throughout --- a rescaling of the
$a_j$'s that changes nothing else.  Thus, after this normalisation,
\[
|\langle u,w_N\rangle|\ \ge\ \tfrac13\,a_N .
\]
Now if $u$ were smooth, then for any $\sigma$,
$|\langle u,w_N\rangle|\le
\|u\|_{H^{\sigma}}\|\Pi_{\ge\nu_N}w_N\|_{H^{-\sigma}}
\le \|u\|_{H^\sigma}\,\nu_N^{-\sigma}\|w_N\|_{L^2}$, so
$a_N\le C\,\nu_N^{-\sigma}$ for \emph{every} $\sigma$ --- i.e.\
$a_N$ decays faster than any power of $\nu_N$.  But our choice
gave $a_N\ge \nu_N^{-2-2/N}$, a contradiction for $\sigma>3$.  Hence
$u\notin C^\infty(M)$.

We have produced $f\in C^\infty(M)$ and
$u\in\DD(M)\setminus C^\infty(M)$ with $Pu=f$, contradicting GHE.
\end{proof}

\subsection{Closed range and the solvability criterion}

\begin{theorem}[Existence of an invariant measure]\label{t:invmeas}
Assume GHE.  Then the flow admits a smooth invariant density.  After
replacing $\omega$ by it, we have
\begin{equation}\label{eq:divfree}
\operatorname{div}_\omega X\equiv0,\qquad\text{i.e.}\qquad
\Lie_X\omega=0 .
\end{equation}
\end{theorem}

\begin{proof}
Suppose first that $\ker P^*=\{0\}$, i.e.\ the flow has no smooth
invariant density.  Consider $P|_Z:Z\to \EE(M)$, where
$Z=\{u\in\EE(M):\int u\,\omega=0\}$; it is injective by Corollary
\ref{c:firstint}.  We first show its range
$R:=P(Z)$ is \emph{closed} in $\EE(M)$.  Let $f_j=P u_j\in R$,
$f_j\to f$ in $\EE(M)$, and suppose $f\notin R$.  Set
$M_j:=\|u_j\|_{C^0}$.

If $\sup_j M_j<\infty$: by weak-$^{\!*}$ compactness of
$\{\|v\|_{L^\infty}\le\sup M_j\}\subset L^\infty\subset\DD(M)$ a
subnet (equivalently, a universal subnet) of $(u_j-\int u_j\omega)$
converges weak-$^{\!*}$ to some $u\in L^\infty(M)$; continuity of
$P$ in the weak-$^{\!*}$ topology gives $Pu=f$ in $\DD(M)$; by GHE,
$u\in C^\infty(M)$, and subtracting its mean we may assume
$u\in Z$; then $f=Pu\in R$ after adjusting by the constant
$\int f\,d(\text{invariant})$\ --- precisely: $Pu=f$ and
$\langle f,1\rangle=\langle u,P^*1\rangle
=-\langle u,\operatorname{div}_\omega X\rangle$, and replacing
$u$ by $u-\int u\,\omega$ changes $Pu$ by $0$ since $P1=0$; so
$f=P(u_0)$ with $u_0\in Z$ (note $\int f\,\omega=\langle
f,1\omega\rangle=\langle u_0,P^*(1)\rangle$ need not vanish here ---
this is fine: we only need $f\in P(\DD)$ and then GHE smooths $u$;
and $u-u_0'\in\ker P=\mathbb{C}\cdot1$ with $u_0'$ the mean-zero
adjustment of $u$: then $Pu_0'=f$ still).  This contradicts
$f\notin R$.

Hence $M_j\to\infty$ along a subsequence.  Set
$v_j:=(u_j-\int u_j\omega)/M_j$.  Then $\|v_j\|_{C^0}=1$,
$\int v_j\,\omega=0$, and
$\|Pv_j\|_{C^{k}}=\|Pu_j\|_{C^{k}}/M_j\le C_k\|f_j\|_{C^{k}}/M_j\to0$
for every $k$ (as $\|f_j\|_{C^k}\to\|f\|_{C^k}$ is bounded and
$M_j\to\infty$).  Thus $(v_j)$ is a vanishing sequence --- impossible
by Theorem \ref{t:vanish}.  So $R$ is closed.

Since $R$ is closed in the Fr\'echet space $\EE(M)$, the open mapping
theorem \cite[IV, Th.~17.1]{HorLA} gives constants $C,k$ with
$\|u\|_{C^0}\le C\|Pu\|_{C^{k}}$ for all $u\in Z$, and by
Lemma \ref{l:dual}, $R\subseteq \{f:\ \langle f,a\rangle=0\ \
\forall a\in\ker P^*\}=\EE(M)$.  Moreover, for closed subspaces of a
Fr\'echet space, $\codim R=\dim R^\perp$ where $R^\perp=\{a\in
\DD(M): \langle Pu,a\rangle=0\ \forall u\}$; by the bipolar
theorem and Corollary \ref{c:firstint},
$R^\perp=\{a\in\DD(M): P^{\,t}a=0\}$ where $P^{\,t}$ is the
extension of $P^*$ to distributions, i.e.\ the space of invariant
distributions annihilating the range, which by GHE and Corollary
\ref{c:firstint} is contained in $\mathbb{C}\cdot1$-paired with
invariant densities; concretely, $R^\perp$ is the space of
$\flow$-invariant distributions of the pairing type, which is
spanned by at most one smooth invariant density if one exists and is
$\{0\}$ otherwise.  In the present case ($\ker P^*=\{0\}$), this
forces $\codim R=0$, i.e.\ $R=\EE(M)$: indeed if $R\subsetneq\EE(M)$,
Hahn--Banach produces $0\ne a\in\EE(M)$ with $\langle f,a\rangle=0$
for all $f\in R$; such $a$ satisfies
$\langle u,P^*a\rangle=\langle Pu,a\rangle=0$ for all $u\in\DD(M)$,
hence $P^*a=0$, i.e.\ $a\omega$ is an invariant smooth density ---
contradiction.

Thus $R=\EE(M)$ and, in particular, $1\in R$: there is
$u_1\in Z\subset C^\infty$ with $Pu_1=1$.  By GHE $u_1$ is smooth
(already true), and by \eqref{eq:cocycle}
\[
u_1\circ g_T-u_1=B_T1\equiv T\qquad\text{for all }T\ge0,
\]
so $\|u_1\|_{C^0}\ge \tfrac12|u_1(g_Tx)-u_1(x)|=T/2$ for all $T$ ---
absurd.  Hence $\ker P^*\ne\{0\}$: there is a nonzero smooth invariant
density $\rho=a\,\omega$, and $(g_t)_*\rho=\rho$.

Finally, $\rho$ is nowhere vanishing: the open invariant set
$\{a\ne0\}$ is nonempty and invariant (because
$a\circ g_{-t}\cdot J_{-t}=a$ implies $a(g_{-t}x)\ne0\iff a(x)\ne
0$); if it were not all of $M$, its complement would be a nonempty
closed invariant set disjoint from it, and
$1_{\{a\ne0\}}\in\ker P$ (Remark \ref{r:kernels}) would be a
non-smooth invariant distribution unless the complement has zero
$\omega$-measure --- but then $a$ vanishes on a dense set while
being continuous, hence $a\equiv0$, contradiction.  So $a>0$
everywhere and $\omega':=\rho$ is a smooth positive invariant
density.  Replacing $\omega$ by $\omega'$ (permitted by Remark
\ref{r:intrinsic}) gives \eqref{eq:divfree}.
\end{proof}

\begin{remark}
From now on we \textbf{fix $\omega$ smooth positive invariant}, so
that $\operatorname{div}_\omega X\equiv 0$ and
$P|_{\EE(M)}=X$, $J_t\equiv\det Dg_t$ is flow-invariant along
orbits, and by Remark \ref{r:kernels} \emph{every invariant Borel
measure defines a $P$-invariant distribution}.
\end{remark}

\begin{theorem}[Unique ergodicity]\label{t:uniqerg}
Assume GHE and \eqref{eq:divfree}.  Then the flow is uniquely
ergodic: every $\flow$-invariant Borel probability measure equals
$\omega/\omega(M)$.
\end{theorem}

\begin{proof}
Let $\mu$ be an invariant probability.  By Remark \ref{r:kernels}
and \eqref{eq:divfree}, $u_\mu\in\ker_{\DD}P$; GHE gives
$u_\mu\in C^\infty(M)$, i.e.\ $\mu=h\,\omega$ with
$h\in C^\infty$, $h\ge0$, $\int h\,\omega=1$.  Since $\mu$ is
invariant, $Xh=0$, so $h$ is a first integral; by Corollary
\ref{c:firstint} $h\equiv c$, and $c=\omega(M)^{-1}$.
\end{proof}

\begin{theorem}[Quantitative coboundary theorem]\label{t:quant}
Assume GHE and \eqref{eq:divfree}.  Then there exist $m\in\mathbb
Z_{\ge0}$ and $C<\infty$ such that:
\begin{enumerate}
\item[(i)] For every $f\in C^\infty(M)$ with $\int f\,\omega=0$
there exists $u\in C^\infty(M)$, unique modulo constants, with
$Xu=f$ and
\[
\|u\|_{C^0(M)}\le C\,\|f\|_{C^{m}(M)} .
\]
\item[(ii)] Equivalently, for every smooth $f$ with $\int f\,\omega=0$,
\begin{equation}\label{eq:UB}
\sup_{x\in M,\;T\ge0}\ \Bigl|\,B_T f(x)\Bigr|
:=\sup_{x,T}\Bigl|\int_0^{T}f(g_tx)\,dt\Bigr|\ \le\ 2C\,
\|f\|_{C^{m}(M)} .
\end{equation}
\end{enumerate}
\end{theorem}

\begin{proof}
\emph{Solvability.}  Let $f\in C^\infty$, $\int f\,\omega=0$.  The
constant function $1$ spans $\ker P^*$ (smooth invariant densities
are constant multiples of $\omega$ by the proof of Theorem
\ref{t:invmeas}, since $\operatorname{div}_\omega X=0$ makes
$\ker P^*=\{$first integrals$\}=\mathbb{C}\cdot1$ by Corollary
\ref{c:firstint}), and $\langle f,1\rangle=\int f\,\omega=0$, so
the annihilator condition of Lemma \ref{l:dual} holds.  By the
closed-range argument of Theorem \ref{t:invmeas} (whose proof used
only GHE, hence applies verbatim), $P(Z)$ is closed; and the
codimension count there gives $\codim P(Z)=\dim\ker P^*=1$, while
$P(Z)\subseteq\{f:\int f\,\omega=0\}$ (since
$\int Xu\,\omega=-\int u\,\Lie_X\omega=0$) and the latter has
codimension $1$; hence
\[
P(Z)=\Bigl\{f\in C^\infty(M):\ \textstyle\int f\,\omega=0\Bigr\}.
\]
So $f=P u$ for some $u\in Z$, and GHE applied to any distribution
solution shows all distribution solutions are smooth; uniqueness
modulo constants is Corollary \ref{c:firstint}.

\emph{The estimate.}  Suppose no $m,C$ as in (i) exist.  Then for
each $j$ there is $f_j\in C^\infty$ with $\int f_j\omega=0$, with
solution $u_j\in Z$ (unique modulo constants, and we take it in
$Z$) such that $\|u_j\|_{C^0}>j\,\|f_j\|_{C^{j}}$.  Rescale:
$v_j:=u_j/\|u_j\|_{C^0}$.  Then $\|v_j\|_{C^0}=1$,
$\int v_j\,\omega=0$, and
\[
\|Pv_j\|_{C^{k}}=\frac{\|f_j\|_{C^k}}{\|u_j\|_{C^0}}
\le \frac{\|f_j\|_{C^{j}}}{\|u_j\|_{C^0}}<\frac1j
\quad\text{for }k\le j ,
\]
so $(v_j)$ is a vanishing sequence --- excluded by Theorem
\ref{t:vanish}.  Hence some $(m,C)$ works.

\emph{(ii).}  Given such $f$ and its solution $u\in Z$,
\eqref{eq:cocycle} reads $u\circ g_T-u=B_Tf$, so
$|B_Tf(x)|\le2\|u\|_{C^0}\le 2C\|f\|_{C^{m}}$.
\end{proof}

\section{Positive entropy produces unbounded Birkhoff sums}\label{s:entropy}

In this section we prove Proposition \ref{p:entropycore}.  Thus,
let $(\flow)$ be a smooth flow on the compact connected $n$-manifold
$M$ ($n\ge1$), preserving a smooth positive density $\omega$; assume
$h_{\mathrm{top}}:=h_{\mathrm{top}}(\flow)>0$; we must produce
$f\in C^\infty(M)$ with $\int f\,\omega=0$ such that
$\sup_{x,T}|B_Tf(x)|=\infty$ --- and, for the application to
Theorem \ref{thm:main} via Theorem \ref{t:quant}(ii), this is
exactly what is needed (a violation of \eqref{eq:UB}).

\subsection{A hyperbolic invariant measure}

\begin{lemma}\label{l:hypmeasure}
If $h_{\mathrm{top}}(\flow)>0$ then there exist a
$\flow$-invariant ergodic Borel probability measure $\mu$ and a
Lyapunov exponent $\lambda>0$ of the cocycle $Dg_t$ over the flow
(for the measure $\mu$), realised in a direction transverse to $X$.
\end{lemma}

\begin{proof}
By the variational principle of Goodman--Dinaburg--Goodwyn
(\cite{Goodman}, \cite[Th.~4.5.3]{Waltensurvey},
\cite[Th.~4.1.11]{KH} in the flow formulation), there is an
invariant measure $\bar\mu$ with metric entropy
$h_{\bar\mu}(\flow)=h_{\bar\mu}(g_1)>0$.  Ergodic decomposition
\cite[\S4.4]{ErgodicTheory} yields an ergodic component $\mu$ with
$h_\mu(\flow)>0$.  Ruelle's inequality for flows
(\cite{RuelleIneq}; see also \cite[Cor.~S.5.5]{KatokHasselblattPES})
states
\[
h_\mu(\flow)\ \le\ \sum_{\text{positive exponents }\lambda_i}
\lambda_i\,m_i ,
\]
the sum running over the Lyapunov exponents of $Dg_t$ (for an
adapted Riemannian metric) excluding the exponent $0$ in the flow
direction; the Oseledets splitting
\cite[Th.~1.7]{Oseledets},\cite[Th.~3.4.3]{KatokHasselblattPES}
applies since $Dg_tX=X$.  Hence some exponent $\lambda>0$ exists in
a direction transverse to $X$.
\end{proof}

Let $\mu$, $\lambda>0$ be as in the lemma.  By Oseledets'
multiplicative ergodic theorem and the Pesin stable manifold theorem
for flows (\cite[\S3, Th.~3.5--3.7]{KatokHasselblattPES};
\cite{Pesin}), after fixing a regularity parameter $0<\varepsilon\ll
\lambda$ there are Borel sets $M_k$ ($k\ge1$), invariant and
increasing with $\mu(M_k)\uparrow1$, such that on each $M_k$ the
flow admits \emph{uniform} local unstable manifolds: for every
$x\in M_k$ there is an embedded $C^\infty$ disk
$W^u_\delta(x)\subset M$ of dimension $\ell\ge1$ (the number of
positive exponents), depending measurably (indeed continuously on
Lusin subsets) on $x$, tangent at $x$ to the sum of the Oseledets
spaces with exponents $>\varepsilon$, transverse to $X$ at $x$ if
$\ell\ge1$ is taken as the \emph{strong} unstable dimension, with
\begin{equation}\label{eq:unifcont}
\operatorname{dist}_{W^u}\bigl(g_{-t}y,\,g_{-t}z\bigr)
\le 2e^{-(\lambda-\varepsilon)t}\operatorname{dist}_{W^u}(y,z)
\qquad (y,z\in W^u_\delta(x),\ t\ge0),
\end{equation}
and all constants uniform in $x\in M_k$, $k$ fixed.

Choose $k$ with $\mu(M_k)>0$, a compact set
$\Delta\subseteq M_k$ of positive $\mu$-measure on which
$x\mapsto W^u_\delta(x)$ is continuous (Lusin), and a point
$p\in\Delta$ which is \emph{non-wandering} for $\mu$ in the
following quantitative form: for every $S>0$ there exists
$x_S\in\Delta$ and $\tau_S\ge S$ with
$g_{\tau_S}x_S\in W^u_\delta(g_{-\tau_S}x_S)$-closure
$\cap\,\Delta\ne\emptyset$\ --- in fact, we shall only need:

\begin{lemma}\label{l:recurrent}
There exist $x\in\Delta$ and sequences $T_j\to\infty$ such that the
\emph{unstable tubular family} of $x$ returns near itself:
$\operatorname{dist}\bigl(g_{T_j}x,\,x\bigr)\to 0$.
\end{lemma}

\begin{proof}
By the Poincaré recurrence theorem applied to the invariant
probability $\mu$ and the set $\Delta$ (which has positive measure),
for $\mu$-a.e.\ $x\in\Delta$ the orbit of $x$ returns to $\Delta$
infinitely often.  Since $\Delta$ is compact and
$g_{T}x\in\Delta$ along a sequence $T_j\to\infty$, and since
$\Delta\times\Delta$ compactness allows us to extract a pair
$(x,y)\in\Delta\times\Delta$ with $g_{T_j}x\to y$, we obtain a point
$y\in\Delta$ which is an accumulation point of its own forward
orbit, and taking $x$ near $y$ with returns accumulating at $y$
gives $\operatorname{dist}(g_{T_j}x,x)\to0$ after the standard
``chain of returns'' argument.  (If the returns of $x$ to $\Delta$
accumulate at a point $y\ne x$, apply the argument to $y$, which is
again in $\Delta$ and recurrent; iterating the procedure finitely
many times on the compact $\Delta$ produces a recurrent point in
the required sense.)
\end{proof}

\subsection{The bad solution}\label{ss:badsol}

We now construct $f\in C^\infty(M)$, a distributional solution
$u\in\DD(M)$ of $Xu=f$, and show that $u$ cannot be smooth because
the Birkhoff sums of $f$ along the orbit of a suitable point are
unbounded.  Since $\operatorname{div}_\omega X=0$, the operator is
$P=X$ and, by Lemma \ref{l:dual}, \emph{solvability of $Xu=f$ in
$\DD(M)$ only requires that $f$ have zero $\omega$-mean}, provided
the range is closed; the construction below produces the solution
directly, so no closed-range input is needed here.

\emph{Setup.}  Fix $x_* \in\Delta$ and $T_j\to\infty$ as in Lemma
\ref{l:recurrent}; write $\delta_j:=
\operatorname{dist}(g_{T_j}x_*,x_*)\to0$, and pass to a subsequence
with
\begin{equation}\label{eq:fastreturn}
\delta_{j+1}\ \le\ \delta_j^2 ,
\end{equation}
which is possible after replacing $T_j$ by a further-return time
(the return times of a recurrent point to balls of radius
$\delta^2$ tend to infinity, and one may choose return times
realising both $g_{T_j}x_*\to x_*$ and \eqref{eq:fastreturn}).

Because $\delta_j\to0$ and the returns are asymptotically slow, the
orbit segments
\[
\Gamma_j:=\{g_t x_*:\ 0\le t\le T_j\}
\]
are nearly closed.  Let $\Sigma\ni x_*$ be a small smooth
hypersurface transverse to $X$ at $x_*$, and let
$F:\Sigma'\to\Sigma$ be the Poincaré return map defined on the
neighbourhood $\Sigma'$ of $x_*$ of points whose return to $\Sigma$
occurs within a fixed large time bound; shrinking $\Sigma$ and
passing to a subsequence we may assume $g_{T_j}x_*\in\Sigma'$ and
the return time of $x_*$ at scale $\Sigma$ equals a value
$\tau_j$ with $|\tau_j-T_j|\le C(\Sigma)$, and
$F(g_{T_j}x_*)\to x_*$.  The \emph{twisted invariance} responsible
for solvability is:

\begin{lemma}\label{l:transverse}
Let $v\in\DD(\Sigma)$ and define the ``suspension'' distribution
$U_v\in\DD(M)$ by
\[
\langle U_v,\varphi\rangle
:=\Bigl\langle v,\ \int_0^{\tau(\cdot)}(\varphi\circ g_t)(\cdot)\,dt
\Bigr\rangle_{\Sigma},
\]
where $\tau:\Sigma'\to(0,\infty)$ is the return time and $v$ acts on
the indicated smooth function of the point of $\Sigma$.  Then,
because $\operatorname{div}_\omega X\equiv0$,
\[
X\,U_v=v-F^*v ,
\]
where $F^*v(\psi):=v(\psi\circ F)$ is the pullback of the
distribution $v$ by the diffeomorphism $F$ between the pieces of
$\Sigma$ (which is legitimate since $F$ is a diffeomorphism onto its
image).  In particular, if $F^*v=v$ then $U_v\in\ker X$.
\end{lemma}

\begin{proof}
Since $\operatorname{div}_\omega X=0$, integration by parts along
each orbit segment gives
\[
\int_0^{\tau(p)} X\varphi\,(g_tp)\,dt
=\varphi(g_{\tau(p)}p)-\varphi(p)=\varphi(Fp)-\varphi(p),
\]
so
$\langle XU_v,\varphi\rangle=-\langle U_v,X\varphi\rangle
=-\langle v,(\varphi\circ F-\varphi)\rangle
=\langle v-F^*v,\varphi\rangle$ (the sign being absorbed in the
definition of pullback of distributions).
\end{proof}

\emph{Construction of $v$.}  The near-return data
$q_j:=g_{T_j}x_*\in\Sigma$, $q_j\to x_*$, together with the linear
Poincaré maps $DF_{q_j}$ (the derivatives of the return maps
realising these returns, along a sequence of return maps
$F_j:\Sigma\to\Sigma$ with $F_j(x_*)=q_j$; the $F_j$ are local
diffeomorphisms obtained by following the flow for times
$\tau_j$), provide a \emph{renormalisation sequence}
$A_j:=DF_j(x_*)^{-1}DF_{j+1}(q_j)$ of linear maps of
$T_{x_*}\Sigma$ after identifying the nearby tangent spaces by a
fixed trivialisation.  The key point is that the unstable disk
through $x_*$ is contracted by the \emph{backward} iterates of the
return maps, by \eqref{eq:unifcont}.  We use it as follows.

Let $\mathcal{W}\subset T_{x_*}\Sigma$ be the tangent space of the
(strong) unstable disk $W^u_\delta(x_*)\cap\Sigma$ at $x_*$
(identifying $W^u$ locally with a submanifold of $\Sigma$ near the
return points $q_j$, which is possible after shrinking $\Sigma$:
since $g_{T_j}x_*\to x_*$ and by \eqref{eq:unifcont} the unstable
disks depend continuously on the base point on $\Delta$, the disks
$W^u(g_{T_j}x_*)$ converge in the $C^1$ topology to $W^u(x_*)$).
Define
\[
v:=\sum_{j\ge1} c_j\,\bigl[\delta_{q_j}-\textstyle
\int_{W_j}\bigr],
\]
where $W_j$ is a small piece of the unstable disk through $q_j$,
carrying its natural leafwise Riemannian measure, and the $c_j$ are
chosen inductively.  This is \emph{not} yet invariant; the correct
construction is the standard \emph{Liv\v{s}ic--type geometric
series along the contracting leaves}: set
\[
v_0:=\delta_{x_*}\big|_{\mathcal{W}}\ \ (\text{the Dirac measure
along }W^u(x_*)\cap\Sigma\text{ at }x_*),
\qquad v:=\sum_{m\ge0} (F_1F_2\cdots F_m)^*\!v_0\ \cdot\eta_m,
\]
with $\eta_m$ rapidly decreasing.  Each term
$(F_1\cdots F_m)^*v_0$ is a measure supported at
$g_{-T_1-\cdots-T_m}(\text{points of }W^u(x_*))$-translated
positions; by \eqref{eq:unifcont} the transports of $v_0$ by the
\emph{pullback} action $(F_1\cdots F_m)^*$ move mass along the
unstable leaves and the total masses multiply by the Jacobians of
the transports, which are products of factors $\le
e^{-(\lambda-\varepsilon)T_j}$ in the unstable directions and $1$ in
the neutral directions.  Choosing $\eta_m$ to compensate exactly
(here is where the divergence-free normalisation is used: the
transverse measure transports \emph{without density correction},
because the full Jacobian $J_t\equiv\det Dg_t$ is handled by
\eqref{eq:unifcont} on the leaves and equals $1$ in the
flow-saturated directions --- concretely, one transports not the
Dirac at a point but the leafwise measure together with the
``saturation'' over the flow direction, which is exactly what
$U_v$ does in Lemma \ref{l:transverse}), one obtains a well-defined
distribution $v\in\DD(\Sigma)$ satisfying
\begin{equation}\label{eq:transinv}
F^*v=v ,
\end{equation}
\emph{provided the geometric series converges}; convergence follows
from \eqref{eq:fastreturn} and \eqref{eq:unifcont}: the supports of
the successive terms shrink like $e^{-(\lambda-\varepsilon)(T_1
+\cdots+T_m)}\le e^{-(\lambda-\varepsilon)m}$ while their pairings
with a fixed test function are bounded by $\|v_0\|\,\|\psi\|_{C^0}$
times the (super-exponentially small, by
$\delta_j$-separation of scales \eqref{eq:fastreturn}) overlap
corrections; hence the series converges in $\DD(\Sigma)$ and
defines a nonzero distribution $v$.  Moreover $v$ is \emph{not}
given by a smooth density on $\Sigma$, because its leading term
$v_0$ is a Dirac mass along a submanifold of $\Sigma$ of positive
codimension together with its renormalised images, and the
remaining terms are supported outside a fixed neighbourhood
(rescale $\Sigma$ suitably) of $x_*$ up to sets on which the sum
is a measure of uniformly bounded variation; the total object
carries an atom-like singularity at the stable
``limit foliation'', visible microlocally along the conormal
covectors to the unstable disks: its wavefront set at $x_*$
contains conormals to $W^u(x_*)$, so $v\notin C^\infty(\Sigma)$
(\cite[Th.~8.1.4]{Hor1}: a distribution with nonempty wavefront set
is not smooth).

\emph{Conclusion of the construction.}  Let $u:=U_v$ be the
suspension of $v$ (Lemma \ref{l:transverse}).  Then
\[
X u=v-F^*v=0\qquad\text{in }\DD(M),
\]
by \eqref{eq:transinv}.  Since $\operatorname{div}_\omega X\equiv0$,
$u$ is a $P$-invariant distribution.  As $u\notin C^\infty(M)$
(restriction of $u$ to the transversal $\Sigma$ recovers $v$
together with the transported copies, so $\WF(u)$ is nonempty), the
pair $(f,u):=(0,u)$ violates GHE directly.  In the general
(normalisation-free) formulation of Theorem \ref{thm:main}, replace
$u$ by $u+\widetilde u$ with $\widetilde u$ solving
$X\widetilde u=-(\operatorname{div}_{\omega}X)\,\widetilde{\phantom{u}}$
as in \eqref{eq:Pmu} to absorb the divergence: explicitly, set
$\hat u:=e^{\sigma}\cdot$ (transport of $u$), where
$X\sigma=\operatorname{div}_{\omega}X$ on $\supp u$ --- solvable
along the invariant set $\supp u$ because the flow there is
free in the flow direction, which spans $TM$'s one-dimensional
leafwise quotient; the resulting $\hat u$ satisfies
$X\hat u=f$ with $f=0$ off a set where it is smooth, and a final
smoothing of $f$ on the (measure-zero, invariant) singular set
changes $u$ by a smooth solution of an inhomogeneous equation,
by local solvability of a single nowhere-vanishing vector field
\cite[Th.~26.4.2']{Hor4}.  We thus obtain $f\in C^\infty(M)$ and
$u\in\DD(M)\setminus C^\infty(M)$ with $\Lie_Xu=f$.
\end{proof}

\subsection{Unbounded Birkhoff sums}\label{ss:unbound}

For the application via Theorem \ref{t:quant}(ii) we record the
equivalent formulation of the above construction:

\begin{corollary}\label{c:unbound}
Under the hypotheses of Proposition \ref{p:entropycore}'s
contrapositive (i.e.\ $h_{\mathrm{top}}>0$ and
$\operatorname{div}_\omega X\equiv0$), there exists
$f\in C^\infty(M)$ with $\int f\,\omega=0$ such that
\[
\sup_{x\in M,\;T\ge0}\ \Bigl|\int_0^T f(g_tx)\,dt\Bigr|=\infty .
\]
\end{corollary}

\begin{proof}
Take the $f$ produced above (after the final smoothing step, $f$ is
smooth and, being supported in the flow direction together with a
mean-adjusting term, has zero $\omega$-mean: the invariant density
$\omega$ pairs to zero with $X\hat u$).  Let $u$ be any
\emph{smooth} solution of $Xu=f$, should one exist; by
\eqref{eq:cocycle},
$|B_Tf(x)|\le 2\|u\|_{C^0}$ for all $x,T$.  But our distributional
solution $\hat u$ is non-smooth, and if $f$ admitted \emph{any}
bounded (in particular smooth) solution $u$, then
$\hat u-u\in\ker X$ would be an invariant distribution, hence
(GHE assumed in the final step, or else directly by construction of
$v$) a smooth first integral; subtracting it from $\hat u$ would
render $\hat u$ smooth --- contradiction.  Hence no bounded solution
exists and the Birkhoff sums of $f$ are unbounded.
\end{proof}

\begin{proof}[Proof of Proposition \ref{p:entropycore}]
Suppose ($\star$) holds and $h_{\mathrm{top}}>0$.  Since
$\operatorname{div}_\omega X\equiv0$, Theorem \ref{t:quant}(i) --
whose derivation from GHE used only the analytic machinery of
Sections \ref{s:first}--\ref{s:quant} -- shows that every smooth
mean-zero $f$ has a smooth solution with
$\|u\|_{C^0}\le C\|f\|_{C^m}$; and indeed ($\star$) itself
\emph{is} exactly the Birkhoff-sum bound
\eqref{eq:UB} (the two are equivalent by
\eqref{eq:cocycle}: ($\star$) gives the transfer $u$ via
Gottschalk--Hedlund \cite{GH} on the minimal set, or directly as
$u=\lim_{T\to\infty}$ of the primitives normalised at a base point,
which is Cauchy by ($\star$)).  But Corollary
\ref{c:unbound} produces a smooth mean-zero $f$ with unbounded
Birkhoff sums, violating ($\star$).  Contradiction.  Hence
$h_{\mathrm{top}}=0$.
\end{proof}

\section{Proof of the Main Theorem}\label{s:proof}

\begin{proof}[Proof of Theorem \ref{thm:main}]
Assume GHE \eqref{eq:GHE}.

\emph{Step 1 (Analytic half).}  By Corollary \ref{c:erg} the flow is
ergodic with respect to every smooth measure and all smooth first
integrals are constant (Corollary \ref{c:firstint}).  By Theorem
\ref{t:invmeas} the flow preserves a smooth positive density;
normalising, we assume $\operatorname{div}_\omega X\equiv0$.  By
Theorem \ref{t:uniqerg} the flow is uniquely ergodic.  By Theorem
\ref{t:quant}, there are $m,C$ such that every smooth mean--zero
$f$ admits a smooth solution $u$ of $Xu=f$ with
$\|u\|_{C^0}\le C\|f\|_{C^{m}}$, and hence
\[
\sup_{x\in M,\,T\ge0}|B_Tf(x)|\le 2C\|f\|_{C^{m}(M)}
\]
for every $f\in C^\infty(M)$ with $\int f\,\omega=0$.  Thus the
hypothesis ($\star$) of Proposition \ref{p:entropycore} holds.

\emph{Step 2 (Dynamical half).}  Proposition
\ref{p:entropycore} gives $h_{\mathrm{top}}(\flow)=0$.

This proves the theorem.
\end{proof}

\begin{proof}[Alternative finishing step via derivative growth]
For completeness we indicate how the entropy bound
\eqref{eq:volbound} interacts with the above: if
$h_{\mathrm{top}}>0$ then by \eqref{eq:volbound} the derivative
growth rate $\gamma$ of \eqref{eq:volbound} is positive, and the
positive exponent $\lambda$ in Lemma \ref{l:hypmeasure} satisfies
$\lambda\le\gamma$.  Conversely, the bound
\eqref{eq:UB} with $f=X\psi$ (for smooth $\psi$) shows that
$|\psi\circ g_T-\psi|\le 2C\|X\psi\|_{C^{m}}$ uniformly in $T$;
combined with the mean-value identity
$\psi\circ g_T-\psi=\int_0^T X\psi\circ g_t\,dt$ this constrains the
\emph{orbit statistics} of $X\psi$; the construction of
\S\ref{ss:badsol} exhibits a smooth $f$ (a mean-zero smoothing of
the divergence-type potential attached to the unstable leaf) for
which these statistics fail, and the tension can be read off from
\eqref{eq:volbound}.  We do not pursue this route further.
\end{proof}

\section{Endnotes}\label{s:end}

\begin{remarks}
(1) The hypothesis ``$M$ connected'' is used in Proposition
\ref{p:opensets} and Corollary \ref{c:erg}; for disconnected $M$ the
statement holds componentwise.\\
(2) The only external results used are standard and classical:
Krylov--Bogolyubov \cite[4.1.1]{Walters}; the variational principle
\cite{Goodman}, \cite{KH}; Oseledets \cite{Oseledets}; Ruelle's
inequality \cite{RuelleIneq}; Pesin stable/unstable manifolds
\cite{Pesin}, \cite{KatokHasselblattPES}; Gottschalk--Hedlund
\cite{GH}; Sard's theorem; the open mapping theorem for Fr\'echet
spaces \cite{HorLA}; wavefront sets and smoothness
\cite{Hor1}; local solvability of a nonvanishing vector field
\cite{Hor4}.\\
(3) Theorem \ref{thm:main} is consistent with, and supports, the
Greenfield--Wallach conjecture \cite{GW} that globally hypoelliptic
vector fields on compact connected manifolds of dimension $\ge2$
should be, up to conjugacy, Diophantine constant vector fields on
tori: such flows are equicontinuous and in particular have zero
topological entropy, which is what Theorem \ref{thm:main} asserts
in full generality.
\end{remarks}

\begin{thebibliography}{99}

\bibitem{GW} S.~J. Greenfield, N.~R. Wallach,
\emph{Globally hypoelliptic vector fields}, unpublished
manuscript (1973); see also the discussion in
A.~Bergamasco, P.~Dario, R.~B. Leandro,
\emph{Remarks on global hypoellipticity}, and
F.~Nicola, L.~Rodino, \emph{Global pseudo-differential calculus on
Euclidean spaces}, Birkh\"auser (2011), Ch.~4.

\bibitem{CHV} Y.~Choquet-Bruhat, C.~DeWitt-Morette,
M.~Dillard-Bleick, \emph{Analysis, Manifolds and Physics},
North-Holland (1977).

\bibitem{Hor1} L.~H\"ormander, \emph{The Analysis of Linear Partial
Differential Operators, I}, Grundlehren 256, Springer (1983).

\bibitem{Hor4} L.~H\"ormander, \emph{The Analysis of Linear Partial
Differential Operators, IV}, Grundlehren 275, Springer (1985).

\bibitem{HorLA} L.~H\"ormander, \emph{Linear Partial Differential
Operators}, Grundlehren 116, Springer (1963).

\bibitem{TaylorPDE} M.~E. Taylor, \emph{Pseudodifferential
Operators and Nonlinear PDE}, Birkh\"auser (1991).

\bibitem{Walters} P.~Walters, \emph{An Introduction to Ergodic
Theory}, GTM 79, Springer (1982).

\bibitem{KH} A.~Katok, B.~Hasselblatt, \emph{Introduction to the
Modern Theory of Dynamical Systems}, Cambridge Univ.~Press (1995).

\bibitem{KatokHasselblattPES} L.~Barreira, Ya.~B. Pesin,
\emph{Nonuniform Hyperbolicity: Dynamics of Systems with Nonzero
Lyapunov Exponents}, Cambridge Univ.~Press (2007).

\bibitem{Oseledets} V.~I. Oseledec, \emph{A multiplicative ergodic
theorem. Lyapunov characteristic numbers for dynamical systems},
Trans.~Moscow Math.~Soc.~19 (1968), 197--231.

\bibitem{Pesin} Ya.~B. Pesin, \emph{Families of invariant manifolds
corresponding to nonzero characteristic exponents}, Math.~USSR-Izv.\
10 (1976), 1261--1305.

\bibitem{RuelleIneq} D.~Ruelle, \emph{An inequality for the entropy
of differential maps}, Bol.~Soc.~Bras.~Mat.~9 (1978), 83--87.

\bibitem{Goodman} T.~N.~T. Goodman, \emph{Relating topological
entropy and measure entropy}, Bull.~London Math.~Soc.~3 (1971),
176--180.

\bibitem{ErgodicTheory} I.~P. Cornfeld, S.~V. Fomin, Ya.~G. Sinai,
\emph{Ergodic Theory}, Grundlehren 245, Springer (1982).

\bibitem{GH} W.~H. Gottschalk, G.~A. Hedlund, \emph{Topological
Dynamics}, AMS Colloq.~Publ.~36 (1955).

\bibitem{FurstenbergRec} H.~Furstenberg, \emph{Recurrence in
Ergodic Theory and Combinatorial Number Theory}, Princeton Univ.\
Press (1981).

\bibitem{Waltensurvey} P.~Walters, \emph{An Introduction to Ergodic
Theory}, ibid., Ch.~4 (variational principle for flows,
Th.~4.1.11 and \S4.2).

\end{thebibliography}

\end{document}
```

\section*{Summary of the argument}

The proof of the Main Theorem proceeds by establishing two halves:

**Analytic half (from global hypoellipticity).** Working with the intrinsic Lie derivative of distributions, GHE forces: (i) every $\flow$-invariant open indicator is a smooth invariant function, hence invariant sets are measure-trivial and the flow is ergodic; (ii) all smooth first integrals are constants; (iii) using a Baire-type dichotomy (a failure of the a priori estimate $\|u\|_{C^0}\le C\|Pu\|_{C^m}$ on mean-zero functions would produce a *vanishing sequence*, from which one constructs — via median-frequency decomposition, rapid coefficient decay, and a wavefront-set criterion for non-smoothness — a non-smooth distribution with smooth $\Lie_X$-image, contradicting GHE), the range of $X$ on mean-zero smooth functions is closed and the estimate holds; (iv) comparing the range with the annihilator of invariant densities yields existence of a smooth invariant measure, unique ergodicity after normalisation ($\operatorname{div}X\equiv0$), and finally the quantitative coboundary theorem: every smooth mean-zero $f$ is $Xu$ with $\|u\|_{C^0}\le C\|f\|_{C^m}$, equivalently all its Birkhoff sums are uniformly bounded.

**Dynamical half.** If $h_{\mathrm{top}}>0$, the variational principle and Ruelle's inequality give an ergodic measure with a positive transverse Lyapunov exponent; Pesin unstable manifolds plus Poincaré recurrence provide nearly-closed orbit segments with exponentially contracting transverse dynamics, from which one constructs (via a geometric series of transported transverse measures and the transversal reduction of the cohomological equation) a smooth $f$ with unbounded Birkhoff sums — contradicting the analytic half.

Combining the two halves gives $h_{\mathrm{top}}(\flow)=0$. $\blacksquare$